\begingroup
\chapter*{Évaluation : géométrie}\markboth{Évaluation : géométrie}{Évaluation : géométrie}\addcontentsline{toc}{chapter}{Évaluation : géométrie}
\renewcommand{\thebmtexo}{DS3.\arabic{bmtexo}}
\renewcommand{\theHbmtexo}{DS3.\arabic{bmtexo}}
\setcounter{bmtexo}{0}
\begin{devoir}[2 h — 20 points]
Calculatrice autorisée. Justifier les réponses et préciser les domaines utiles. Les questions sont indépendantes sauf indication. 
\end{devoir}
\exo[Produit scalaire et barycentre]\label{t11s-ds3-e01}\bareme{6 points}
Dans un repère orthonormé, $A(0;0),B(6;0),C(2;4)$.
\begin{enumerate}\item Calculer $\vect{AB}\cdot\vect{AC}$ et l’aire du triangle (2 points).
\item Calculer le barycentre $G$ de $(A,1),(B,1),(C,2)$ (2 points).
\item Déterminer le lieu $MA^2+MB^2=26$ (2 points).\end{enumerate}
\begin{corrige}[exercice~\ref{t11s-ds3-e01}]\label{sol-t11s-ds3-e01}
1. Produit $12$; aire $6\times4/2=12$ (1+1).
2. Somme $4$, $G=((0+6+4)/4;(0+0+8)/4)=(5/2;2)$ (condition 1, coordonnées 1).
3. $I=(3;0)$, somme $2MI^2+18=26$, donc cercle de centre $I$ et rayon $2$ (réduction 1, lieu 1).
\end{corrige}
\exo[Trigonométrie]\label{t11s-ds3-e02}\bareme{4 points}
\begin{enumerate}\item Déterminer la mesure principale de $19\pi/6$ (1 point).
\item Résoudre $2\cos^2x-1=0$ sur $[0;2\pi[$ (2 points).
\item Résoudre $\sin x\ge1/2$ sur le même intervalle (1 point).\end{enumerate}
\begin{corrige}[exercice~\ref{t11s-ds3-e02}]\label{sol-t11s-ds3-e02}
1. $19\pi/6-4\pi=-5\pi/6$.
2. $\cos2x=0$, donc $x=\pi/4+k\pi/2$, soit $\pi/4,3\pi/4,5\pi/4,7\pi/4$ (famille 1, intervalle 1).
3. $[\pi/6;5\pi/6]$.
\end{corrige}
\exo[Similitude]\label{t11s-ds3-e03}\bareme{4 points}
Étudier $x'=-2y+1$, $y'=2x+2$ : rapport et angle, centre, image de $(1;0)$ et facteur d’évolution des aires.
\begin{corrige}[exercice~\ref{t11s-ds3-e03}]\label{sol-t11s-ds3-e03}
Rapport $2$, angle $\pi/2$ (1). Point fixe $x=-2y+1$, $y=2x+2$, donnant $x=-3/5,y=4/5$ (1). Image $(1;4)$ (1); facteur des aires $4$ (1).
\end{corrige}
\exo[Un plan dans l’espace]\label{t11s-ds3-e04}\bareme{6 points}
En repère orthonormé direct, $A(1;0;0),B(0;1;0),C(0;0;1)$.
\begin{enumerate}\item Calculer $\vect{AB}\wedge\vect{AC}$ et donner une équation de $(ABC)$ (2 points).
\item Déterminer l’intersection avec la droite $(x;y;z)=(t;2t;3t)$ (2 points).
\item Calculer la distance de l’origine au plan et l’aire du triangle $ABC$ (2 points).\end{enumerate}
\begin{corrige}[exercice~\ref{t11s-ds3-e04}]\label{sol-t11s-ds3-e04}
1. Vectoriel $(1;1;1)$; plan $x+y+z=1$ (1+1).
2. $6t=1$, donc $(1/6;1/3;1/2)$, unique (1+1).
3. Distance $1/\sqrt3$; aire $\sqrt3/2$ (1+1).
\end{corrige}
\endgroup
